\section*{Chapter \ref{ch:b2:prices-as-information} --- Prices as Information}

\begin{solution}{exo:b2:prices-as-information:1}
% ledger: b2:sw2010-returns
% figdata: returns_summary
A dollar on the favourite returns 94.5 cents on average ($1 - 0.055$); a dollar at odds of 100/1 or longer
returns 39 cents ($1 - 0.61$). If the odds were unbiased forecasts, every bar would have the same length:
every kind of bet would lose the track's share of the pool and nothing more.
\end{solution}

\begin{solution}{exo:b2:prices-as-information:2}
% code: exercises/solutions.py counts
Each household's 40 budget shares sum to one, so 39 numbers describe it: $2{,}000{,}000 \times 39 = 78$
million numbers. A market posts 39 prices. The planner would also need every firm's technology and every
household's \omterm{def:b2:exchange-economies:endowment}{endowment}, and a way to make people report truthfully.
\end{solution}

\begin{solution}{exo:b2:prices-as-information:3}
% code: exercises/solutions.py two_signal
The naive price is $(2 - 0.5)/2 = 0.75$; the \omterm{def:b2:prices-as-information:fully-revealing}{fully revealing} price is $2 - 0.5 = 1.5$. At 0.75, a trader who
learnt $s_1 + s_2 = 1.5$ would expect the asset to be worth 1.5 and would want to buy, whichever group she
belongs to: the naive price is not an equilibrium once traders learn from it.
\end{solution}

\begin{solution}{exo:b2:prices-as-information:4}
Prices that are badly calibrated across many contracts (contracts priced near 0.8 whose events happen far
less or far more than 80 per cent of the time); forecast errors that public information available at the
time (polls, earlier prices) could have predicted; or prices that do not move when traders learn news
that is relevant to the event.
\end{solution}

\begin{solution}{exo:b2:prices-as-information:5}
% code: exercises/solutions.py gs
$k = e^{0.1} - 1 = 0.1052$. Then $\phi^* = 0.25\sqrt{(1 - 0.1052)/0.1052} = 0.25 \times 2.917 = 0.729$, and the
error variance of the price's forecast is $k = 0.105$.
\end{solution}

\begin{solution}{exo:b2:prices-as-information:6}
% code: exercises/solutions.py gs
A dearer signal should attract fewer buyers and leave the price less informative. At $c = 0.2$, $k = e^{0.4} -
1 = 0.492$ and $\phi^* = 0.25\sqrt{0.508/0.492} = 0.254$: the informed share falls from 0.433 to 0.254 and the
error variance rises from 0.25 to 0.492. If $\sigma_z$ doubled instead, $\phi^*$ would double, to 0.866 at
$c = 0.1116$, and the error variance would stay at 0.25: more noise attracts exactly enough \omterm{def:b2:prices-as-information:informed-trader}{informed
traders} to offset it.
\end{solution}

\begin{solution}{exo:b2:prices-as-information:7}
% ledger: b2:bfnr2008-polls
% code: exercises/solutions.py iem luck
$(1.91 - 1.49)/1.91 = 22$ per cent smaller with election-eve prices, $(1.91 - 1.58)/1.91 = 17$ per cent with
last-week prices. The market was closer in $9/15 = 60$ per cent of the elections. If each election were a
fair coin flip between market and polls, at least 9 wins out of 15 would happen with probability 0.304; so
9 of 15 is weak evidence on its own. One needs the size of the errors in each election, how independent the
elections are, and the number of elections, which only a longer record supplies.
\end{solution}

\begin{solution}{exo:b2:prices-as-information:8}
% code: exercises/solutions.py payback
A return of $-5.5$ per cent means the average dollar comes back as 0.945, so these horses win $0.945/3 =
0.315$ of their races; the odds suggest $1/3 = 0.333$. They win $0.315/0.333 = 94.5$ per cent as often as the
odds say. The horse at 100/1 wins $0.0039/0.0099 = 39$ per cent as often as its odds say. Both shortfalls
include the track's share; the difference between them is the bias.
\end{solution}

\begin{solution}{exo:b2:prices-as-information:9}
Follow the price over the next hours and days, and compare it with any other market or poll on the same
election. If, without news, the price falls back towards 0.40 as other traders sell, the push was treated
as noise and did not last: the market resisted it. If the price stays near 0.55 while other forecasts do
not move, either the other traders believed the buyer had information, or the market is too thin for
them to correct it; the volume of their selling helps tell the two apart. If later news justifies 0.55, the
buyer was informed, not manipulating.
\end{solution}

\begin{solution}{exo:b2:prices-as-information:10}
(a) Wealth is $W = x(v - p)$, normal with mean $x(m - p)$ and variance $x^2\sigma^2$. Then $\mathbb{E}[-e^{-aW}] =
-\exp(-ax(m - p) + a^2x^2\sigma^2/2)$, which is largest where the exponent is smallest: $-a(m - p) + a^2x\sigma^2 =
0$, so $x = (m - p)/(a\sigma^2)$. (b) If $p = s_1 + s_2$, every trader knows $s_1 + s_2$ and expects $v$ to be $s_1 + s_2$
with variance $\sigma_\varepsilon^2$; her demand is $(s_1 + s_2 - p)/(a\sigma_\varepsilon^2) = 0$. Demands sum to
zero, so the market clears, and the price reveals $s_1 + s_2$, all the information traders hold together. The
variances of the signals do not matter, because once $s_1 + s_2$ is known each signal adds nothing.
\end{solution}

\begin{solution}{exo:b2:prices-as-information:11}
With $\sigma_z = 0$ and $\phi > 0$, the price reveals $\xi = \theta$ exactly, so an uninformed trader knows as much as
an informed one without paying; each \omterm{def:b2:prices-as-information:informed-trader}{informed trader} would do better by not buying, so $\phi > 0$ is not an
equilibrium. With $\phi = 0$ the price carries no information about $\theta$, and $\mathrm{Var}(v \mid p) = 2$.
A trader who bought the signal would have the utility ratio $e^{c}\sqrt{1/2}$, which is below one, so better,
when $c < \tfrac12\ln 2$. Then someone would buy, and $\phi = 0$ is not an equilibrium either.
\end{solution}

\begin{solution}{exo:b2:prices-as-information:12}
% code: exercises/solutions.py spending
Spending per trader is $\phi^* c = 0.433 \times 0.1116 = 0.0483$. The asset pays $\theta + \varepsilon$ whoever knows
$\theta$, so trading only moves wealth between traders: the informed gain from the noise traders whose
supply they absorb at better prices, and the uninformed gain from reading the price. In this model the
spending is a social cost and the gains are transfers. Better prices have real value when decisions with
real consequences depend on them: how much to plant, store or import, which firms to fund. The model leaves
those out.
\end{solution}

\begin{solution}{pb:b2:prices-as-information:1}
% code: exercises/solutions.py harvest gs record expected_msq seasons_needed
\textbf{Part I.}
1. $16 + 16 + 9 = 41$.
2. With both signals the error is $\varepsilon$, variance 9; with one, $16 + 9 = 25$; with none, 41.
3. Two numbers a year, but each is spread among the region's farmers: the planner must collect every farmer's
observation and combine them.
4. If a low forecast brings cheaper seed or more aid, northern farmers gain by reporting a low $s_N$: the
planner's information depends on reports that people have reasons to bend.
5. 9.
6. \Cref{prop:b2:prices-as-information:revealing}: a \omterm{def:b2:prices-as-information:fully-revealing}{fully revealing} price shows $s_N + s_S$ to everyone.
\textbf{Part II.}
7. Northern traders expect $106$ with variance 25 and demand $(106 - p)/25$; southern traders $(98 - p)/25$. The
sum is zero at $p = 102$.
8. $(106 - 102)/25 = 0.16$ contracts.
9. $p = 100 + 6 - 2 = 104$.
10. Every trader then expects 104 with variance 9 and demands $(104 - 104)/9 = 0$: the market clears.
11. The price 104, with error variance 9. The naive price has error variance $\mathrm{Var}((s_N + s_S)/2 + \varepsilon)
= 8 + 9 = 17$.
12. $s_N = 104 - 100 - s_S = 6$: the northern farmers' signal.
\textbf{Part III.}
13. $k = e^{0.16} - 1 = 0.1735$.
14. $\phi^* = 0.25\sqrt{0.8265/0.1735} = 0.25 \times 2.183 = 0.546$.
15. $\mathrm{Var}(\theta \mid p) = k = 0.174$.
16. At $c = 0.05$ the informed share rises to 0.729 and the error variance falls to 0.105.
17. If it did, no trader would pay 0.08 for a forecast that the price gives free, and then the price would
contain no forecast: the \omterm{def:b2:prices-as-information:gs-paradox}{Grossman--Stiglitz paradox}.
18. The \omterm{def:b2:prices-as-information:informed-trader}{informed traders} pay, and recover the cost by trading against the noise in supply; the uninformed
read the price free.
\textbf{Part IV.}
19. Office: $9 + 4 = 13$; market: $9 + 2.25 = 11.25$.
20. The office's errors are $-2.3, -8.2, -7.4, 3.7, 1.4, 0.2, -1.9, -3.2, 2.9, 2.1, 0.3, 3.2$, the market's
$1.6, -5.0, -6.7, 4.1, 0.7, -5.9, 0.4, -1.3, 1.7, -3.7, 3.4, 4.1$. Mean absolute errors: $36.8/12 = 3.07$ for the
office, $38.6/12 = 3.22$ for the market. Mean squared errors: $179.98/12 = 15.00$ and $171.36/12 = 14.28$.
21. In 7 of the 12 seasons.
22. The standard error of the mean difference is $17.8/\sqrt{12} = 5.1$; $-0.72$ is far less than one standard
error from zero. The record cannot tell the forecasts apart. Even the expected difference, $11.25 - 13 = -1.75$,
would be two standard errors from zero only after $(2 \times 17.8/1.75)^2 = 414$ seasons.
23. The average's errors sum to 33.4 in absolute value: a mean absolute error of $33.4/12 = 2.78$. Its expected
squared error is $9 + (4 + 2.25)/4 = 10.56$, below both forecasts, because the two errors are independent.
24. With a forecast costing 0.08, 54.6 per cent of traders buy it and the price forecasts $\theta$ with error
variance 0.174. Over twelve seasons the market and the office cannot be told apart: most of each error is the
weather that no one forecasts. The model gives the market a small edge (11.25 against 13) and the average of the
two a larger one (10.56).
\end{solution}
